\chapter{모스 부호}
% full version: chapters/04_code.tex

모스 부호에는 점과 선, 두 기호와 그 사이의 쉼이 있다. 뉴런의 알파벳은 그보다도 빈약하다. 기호는 짧은 딸깍 하나뿐이다. 그렇다면 메시지 전체는 쉼 속에 적혀 있다는 뜻이다. 뉴런은 대체 어떤 언어로 이야기하는가?

\section*{딸깍으로 얼마나 말할 수 있나}

통신 엔지니어처럼 시작하자. 이런 채널에는 정보가 얼마나 들어가는가? 수신자가 딸깍의 도착 시각을 1밀리초 정밀도로 구별한다면 딸깍 하나가 8비트 정도를 나를 수 있다. 수신자가 1초 동안 딸깍이 몇 번 왔는지만 센다면, 같은 흐름에서 얻는 정보는 수십 배 적다. 채널 용량은 거의 모두 \emph{시간} 속에 있다.

\pencilfig{4}{\pencilart{4}}{위는 모스 부호, 아래는 뉴런이다. 기호는 딸깍 하나뿐이고, 메시지 전체는 쉼 속에 적혀 있다.}

그런데 채널에는 잡음이 많다. 뜻밖에도 뉴런 자체는 정확하다. 똑같이 변하는 전류를 여러 번 넣으면 뉴런은 1밀리초 정밀도로 매번 같은 순간에 딸깍한다. 잡음은 연결에서 생긴다. 전선 끝까지 달려간 딸깍의 절반 넘게가 그냥 사라진다. 물질이 매번 방출되지는 않기 때문이다. 살아 있는 피질에서 딸깍의 흐름은 거의 무작위로 보인다. 이것은 잡음인가, 아니면 우리가 아직 읽지 못하는 메시지인가? 잘 압축된 메시지는 제3자에게 무작위와 구별되지 않으므로, 이 물음은 아직 열려 있다.

\section*{느리지만 확실하게}

가장 단순한 부호는 딸깍을 세는 것이다. 자주 올수록 큰 수이다. 이 부호는 손실에도 떨림에도 끄떡없다. 하지만 느리다. 빈도를 10퍼센트 정밀도로 알려면 딸깍이 백 번 필요하고, 그러려면 몇 초가 걸린다. 그런데 여러분은 한순간 비친 그림 속 동물을 0.15초쯤이면 알아보고, 그동안 신호는 처리 단계를 열 개쯤 지난다. 단계마다 각 뉴런은 기껏해야 한 번 딸깍할 틈밖에 없다.

해법은 뉴런을 여럿 한꺼번에 쓰는 것이다. 청취자 한 명이 1분 동안 딸깍을 세는 대신 청취자 백 명이 한순간에 센다. 그런데 여기에도 함정이 있다. 이웃한 뉴런들의 잡음은 군중 속 소문처럼 서로 조금씩 닮았다. 이웃 수십 개로 평균하면 도움이 되지만, 그 이상은 거의 효과가 없다.

\pencilfig{122}{%
% error of the group-average rate relative to one neuron, sqrt(c + (1-c)/N): c = 0.1 (neighbours' noise
% is slightly alike; full edition, chapter 4) and c = 0 (independent noise). x = 0.8 log2 N, y = 2.2 * error
\draw[pen] (0,0) -- (5.9,0); \draw[pen] (0,0) -- (0,2.45);
\foreach \x/\n in {0/1, 2.66/10, 5.31/100}{\draw[pcGray, line width=0.5pt] (\x,0) -- (\x,-0.08); \node[lbl, anchor=north] at (\x,-0.08) {\n};}
\node[lbl, anchor=north] at (2.95,-0.38) {평균하는 뉴런 수};
\node[lbl, anchor=south, rotate=90] at (-0.1,1.2) {오차};
\draw[pcGray, densely dashed, line width=0.5pt] (0,0.70) -- (5.6,0.70);
\draw[penlite=pcBlue] plot[smooth] coordinates {(0.00,2.20) (0.47,1.80) (0.80,1.56) (1.27,1.27) (1.60,1.10) (2.07,0.90) (2.40,0.78) (2.87,0.64) (3.20,0.55) (3.67,0.45) (4.00,0.39) (4.47,0.32) (4.80,0.28) (5.27,0.22) (5.60,0.19)};
\draw[penlite=pcRed] plot[smooth] coordinates {(0.00,2.20) (0.47,1.84) (0.80,1.63) (1.27,1.39) (1.60,1.25) (2.07,1.10) (2.40,1.01) (2.87,0.92) (3.20,0.87) (3.67,0.82) (4.00,0.79) (4.47,0.76) (4.80,0.74) (5.27,0.73) (5.60,0.72)};
\node[lbl, text=pcRed, anchor=west, align=left] at (5.7,0.74) {닮은 잡음\\(피질처럼)};
\node[lbl, text=pcBlue, anchor=west] at (5.7,0.19) {독립 잡음};
\node[lbl, anchor=west] at (0.9,2.15) {뉴런 하나};
\node[lbl, anchor=north west] at (0.05,0.66) {3분의 1};
}{뉴런 무리로 평균한 빈도의 오차. 잡음이 독립이라면 오차는 계속 줄어들 것이다. 이웃의 닮은 잡음은 오차를 약 3분의 1에서 멈추게 하고, 이 한계에는 뉴런 수십 개면 이미 도달한다.}

\section*{빠르게: 누가 먼저인가}

두 번째 해법은 수를 시간에 적는 것이다. 센 신호는 뉴런을 일찍 딸깍하게 하고, 약한 신호는 늦게 딸깍하게 한다. 딸깍 단 하나가 이미 수를 나른다. 그런데 시계가 없다면 시간을 어디서부터 재는가? 뉴런끼리 비교하면 된다. 누가 먼저 딸깍했고 누가 두 번째였는가. 뉴런 열 개의 딸깍 도착 순서는 20비트가 넘는 정보를 나른다. 게다가 모두가 함께 내는 일제 발화가 시작 표시 역할을 한다. 모스 부호에서 단어 사이의 긴 쉼과 같다.

\pencilfig{121}{%
% latency code: one click per neuron; the stronger the input, the earlier the click
\foreach \y/\t/\l/\k in {1.5/1.1/{센 입력}/1, 0.95/2.4/{중간}/2, 0.4/4.2/{약한 입력}/3}{
  \draw[pen] (0.4,\y) -- (5.1,\y);
  \draw[pcBlue, line width=0.9pt] (\t,\y) -- (\t,\y+0.42);
  \node[lbl, anchor=east] at (0.3,\y+0.1) {\l};
  \draw[dotpen=pcAmber, wash=pcAmber] (5.6,\y+0.16) circle (0.17);
  \node[font=\scriptsize, text=pcAmber!60!black] at (5.6,\y+0.16) {\k};}
\draw[pcGray, densely dashed, line width=0.5pt] (0.55,0.25) -- (0.55,2.1);
\node[lbl, anchor=south] at (0.55,2.1) {자극};
\node[lbl, text=pcAmber!70!black, anchor=south] at (5.6,2.1) {순서};
\draw[arr] (0.7,0.05) -- (4.9,0.05); \node[lbl, anchor=north] at (2.8,0.02) {시간};
}{센 입력은 이른 딸깍, 약한 입력은 늦은 딸깍. 시작 순간을 몰라도 딸깍의 순서가 어느 입력이 더 센지 알려 준다.}

\section*{부호는 듣는 쪽이 고른다}

같은 딸깍 열이 빈도도, 시각도, 순서도 함께 나른다. 그중 무엇을 읽을지는 수신자가 정한다. 양동이 “구멍이 넓은” 뉴런은 물을 천천히 모으므로 빈도만 듣는다. 빨리 잊는 뉴런은 동시성만 듣는다. 그리고 앞에서 보았듯이 브레이크는 뉴런을 한 모드에서 다른 모드로 바꿀 수 있다.

연결조차 신호를 거른다. 어떤 연결은 금방 지쳐서 빈도가 아니라 빈도의 \emph{변화}를 알린다. 반대로 어떤 연결은 연달아 오는 몇 번의 딸깍, 즉 버스트를 더 잘 통과시킨다. 뉴런에게 “선”이 생기는 것이다. 그사이 억제 뉴런은 문장 부호를 찍는다. 흐름을 창으로 자르고, 너무 시끄러우면 소리를 줄인다.

\section*{아껴 쓰는 언어}

마지막으로, 이 언어는 비싸다. 딸깍 하나하나에 에너지가 들고, 뇌 전체는 약 20와트를 쓴다. 이 전력을 모든 뉴런에 나누면 뉴런 하나당 평균 1초에 한 번 정도의 딸깍이 된다. 대부분의 뉴런은 대부분의 시간 동안 조용하다. 뇌의 부호는 인색할 수밖에 없다.

모스 부호에서 영어에서 가장 흔한 글자 E는 점 하나이다. 흔한 것은 짧게 부호화한다. 뉴런도 비슷하다. 예상된 것에는 딸깍을 적게 쓴다. 변하지 않는 자극은 점점 응답을 일으키지 않게 되고, 감지기는 변화를 알리며, 도파민은 뜻밖의 일만 알린다. 뇌는 딸깍을 새 소식에 쓴다.

\fullversion{4}
